\section{Conclusion}\label{sec:conclusion}

We proposed a weighting and an offset for graph collaborative filtering under exposure bias. Sampled BPR
negatives are down-weighted by one minus the probability that an unlogged pair is positive, and the standardised
logit of an inverse-propensity-weighted additive imputation enters the score as a fixed offset. Both nuisances
are fitted once on the biased log. Uniformly collected data are used only for validation and testing, except that
the shipped Coat propensities were estimated by the dataset authors from a uniform sample.

On Coat, Yahoo! R3, KuaiRand and KuaiRec the final method (\qp) has a higher mean than every graph baseline we
ran, and the differences are significant except against LightGCN-pt on KuaiRec. It is the best method we ran on
Yahoo! R3 and KuaiRand. In the main run it is above iALS on Coat and below DR-JL and MRDR on KuaiRec. Across retraining
seeds these differences are not significant. Training-free doubly robust propagation remains the best method we
ran (0.6378). The ablations show that the gain depends on an informative posterior: a shuffled control removes it
on Coat and Yahoo! R3. On Coat (unmatched grids) and KuaiRec, the offset alone recovers most of the improvement.

Four extensions follow from the limitations. Computing the posterior on the fly would remove the memory limit of
dense matrices on large catalogues. Matched offset-only runs on Yahoo! R3 and KuaiRand would test the offset claim
beyond Coat. A faithful reproduction of AdvInfoNCE, with its original adversary, would replace our simplified
re-implementation. Propensities estimated without the uniform data used for testing would remove the possible
leak on Coat.
